\documentclass{article}
\input{../../../lib.sty}
\title{\texttt{fn sample\_truncated\_negative\_binomial\_rational}}
\author{Michael Shoemate}
\begin{document}
\maketitle

This document proves that the implementation of \texttt{sample\_truncated\_negative\_binomial\_rational}
in \texttt{mod.rs} satisfies its proof definition.

\section{Hoare Triple}
\subsection*{Precondition}
$\texttt{eta} \in \mathbb{Q}_{>0}$ and $\texttt{x} \in \mathbb{Q}_{>0}$.

\subsection*{Pseudocode}
\lstinputlisting[language=Python,firstline=2,escapechar=|]{./pseudocode/sample_truncated_negative_binomial_rational.py}

\subsection*{Postcondition}
\label{postcondition}
For any setting of the input parameters \texttt{eta} and \texttt{x} such that the given preconditions hold,
\texttt{sample\_truncated\_negative\_binomial\_rational} either returns \texttt{Err(e)} due to a lack of system entropy,
or \texttt{Ok(out)}, where \texttt{out} is distributed as a negative binomial random variable
with shape \texttt{eta} and success probability $1 - e^{-x}$, parameterized as the number of failures before
the \texttt{eta}-th success, conditioned on \texttt{out > 0}.

\section{Proof}
Assume the preconditions and put $q=e^{-x}$ and $p=1-q$, so $0<p,q<1$.
For $k\ge1$, the target distribution is
\begin{equation}
\label{eq:target}
\Pr[\texttt{out}=k]
=\frac{q^k}{p^{-\eta}-1}\frac{(\eta)_k}{k!},
\qquad (\eta)_k:=\prod_{j=0}^{k-1}(\eta+j).
\end{equation}

\begin{lemma}
\label{lem:geometric}
For $\eta=1$, the output follows the target law.
\end{lemma}
\begin{proof}
Line \ref{line:geometric} returns one plus
\rustdoc{traits/samplers/cks20/fn}{sample\_geometric\_exp\_fast}.
By its proof definition, this is a sample on the positive integers with
$\Pr[\texttt{out}=k]=p q^{k-1}$, which agrees with
Equation~\ref{eq:target} when $\eta=1$.
\end{proof}

\begin{lemma}
\label{lem:fractional}
For $0<\eta<1$, the output follows the target law.
\end{lemma}
\begin{proof}
Let $Y=G+1$, where $G$ is the geometric draw on line
\ref{line:positive_draw}. By its proof definition,
$\Pr[Y=k]=p q^{k-1}$ for $k\ge1$.

The call to \rustdoc{traits/samplers/negative_binomial/fn}{sample\_integer\_envelope\_accept}
on line \ref{line:positive_accept} uses
$\texttt{eta\_numer}/\texttt{eta\_denom}=\eta+1$ and
$\texttt{base\_denom}/\texttt{eta\_denom}=2$.
Consequently, conditional on $Y=k$, its acceptance probability is
\begin{equation*}
A(k)=\prod_{i=0}^{k-2}\frac{\eta+1+i}{2+i}
=\prod_{i=1}^{k-1}\frac{\eta+i}{i+1}.
\end{equation*}
Each factor belongs to $(0,1]$. In particular, $A(1)=1$ by the empty-product
convention. The accepted joint mass is
\begin{align*}
\Pr[Y=k,\text{ accepted}]
&=p q^{k-1} A(k)
=\frac{p}{\eta q}\,q^k\frac{(\eta)_k}{k!}.
\end{align*}
The total acceptance probability is positive (already the event $Y=1$
is accepted with probability $p>0$), so the rejection loop terminates
almost surely. By the binomial series,
\[
\sum_{k=1}^{\infty}q^k\frac{(\eta)_k}{k!}
=(1-q)^{-\eta}-1=p^{-\eta}-1.
\]
Conditioning the displayed joint mass on acceptance therefore yields
Equation~\ref{eq:target}.
\end{proof}

\begin{lemma}
\label{lem:integer-envelope}
For $\eta>1$, the output follows the target law.
\end{lemma}
\begin{proof}
Every iteration draws an independent ordinary negative-binomial variable
$K$ via line \ref{line:ordinary_draw} and accepts it when $K>0$ on
line \ref{line:positive_branch}.
By the proof definition of the ordinary sampler,
\[
\Pr[K=k]=p^\eta q^k\frac{(\eta)_k}{k!}
\quad(k\ge0).
\]
Since $\Pr[K>0]=1-p^\eta>0$, the loop terminates almost surely, and
conditioning on positivity gives
\[
\Pr[\texttt{out}=k]
=\frac{p^\eta q^k(\eta)_k/k!}{1-p^\eta}
=\frac{q^k(\eta)_k/k!}{p^{-\eta}-1}
\quad(k\ge1).
\]
\end{proof}

\begin{proof}[Proof of the postcondition]
The three lemmas cover every valid $\eta>0$ and establish the exact
distribution. The only fallible operations invoke
\rustdoc{traits/samplers/cks20/fn}{sample\_geometric\_exp\_fast},
\rustdoc{traits/samplers/negative_binomial/fn}{sample\_integer\_envelope\_accept},
or \rustdoc{traits/samplers/negative_binomial/fn}{sample\_negative\_binomial\_rational}.
Their proof definitions restrict returned errors to lack of system entropy;
other arithmetic operations are exact on arbitrary-precision integers and
rationals. Hence the function satisfies the postcondition.
\end{proof}

\end{document}
